\chapter{Modelle organisieren}\label{kap:organisieren}
\index{Organisieren}

Ein optisches Modell besteht schnell aus Dutzenden oder Hunderten von Objekten: Linsenflächen,
Fassungen, Blenden, Hilfskurven, eingelesene Baugruppen. Dieses Kapitel zeigt die Werkzeuge,
mit denen Sie Ordnung halten -- Ebenen, Gruppen und den Modellbaum --, das Dock
\gui{Properties} mit den Eigenschaften und freien Attributen eines Objekts und schließlich das
Rückgängigmachen.

Ebenen und Gruppen ergänzen sich:

\begin{itemize}
  \item \textbf{Ebenen} sind \emph{flach} und ordnen nach \emph{Art}: Hilfsgeometrie,
    Gehäuse, Optik, Strahlen als Kurven. Jedes Objekt liegt auf genau einer Ebene; die Ebene
    bestimmt seine Farbe.
  \item \textbf{Gruppen} sind \emph{verschachtelt} und ordnen nach \emph{Zugehörigkeit}: ein
    Objektiv mit seinen Linsen, eine Baugruppe aus einer STEP-Datei. Jedes Objekt liegt in
    höchstens einer Gruppe.
\end{itemize}

Beide haben einen eigenen Sichtbarkeitsschalter; ein Objekt ist nur zu sehen, wenn es selbst,
seine Ebene und alle seine Gruppen eingeschaltet sind (\autoref{sec:ausblenden}).

\section{Ebenen}\label{sec:ebenen}
\index{Ebene}\index{Layer}

Die Ebenen verwaltet das Dock \gui{Layers} (\menu{View,Layers}). Es zeigt eine Liste mit den
Spalten \gui{Layer} (Kontrollkästchen und Name), \gui{Lock} (Kontrollkästchen) und
\gui{Colour} (Farbfeld) und darunter vier Knöpfe. Die \emph{aktuelle} Ebene steht fett. Ein
Rechtsklick auf eine Ebene öffnet ein Kontextmenü mit denselben Befehlen (\gui{Make current},
\gui{Rename}, \gui{Colour...}, \gui{Lock}/\gui{Unlock}, \gui{Move selection here},
\gui{Delete layer...}).

\begin{referenz}{>{\raggedright\arraybackslash}p{4.2cm}X}
\toprule
\textbf{Bedienung} & \textbf{Wirkung} \\
\midrule
\endhead
Kontrollkästchen & Ebene ein- oder ausblenden. Objekte einer ausgeblendeten Ebene sind
  unsichtbar und lassen sich weder auswählen noch fangen. \\
\maus{Doppelklick} auf den Namen & Name bearbeiten; mit \keys{Enter} bestätigen
  (ein Rückgängig-Schritt \enquote{Rename layer}). \\
Kontrollkästchen \gui{Lock} & Ebene sperren oder entsperren. Objekte einer gesperrten Ebene
  bleiben sichtbar und lassen sich fangen, aber nicht auswählen -- und damit auch nicht
  verschieben, bearbeiten oder löschen. Beim Sperren verlassen sie die Auswahl. \\
\maus{Doppelklick} auf das Farbfeld & öffnet den Dialog \gui{Layer colour}. Die neue Farbe
  gilt sofort für alle Objekte der Ebene, die keine eigene Farbe haben (ein Rückgängig-Schritt
  \enquote{Layer colour}). \\
\gui{New layer} & legt eine neue Ebene \enquote{Layer \textit{n}} an, mit einer
  automatisch gewählten Farbe, die sich von den bisherigen unterscheidet (ein
  Rückgängig-Schritt \enquote{New layer}). \\
\gui{Delete layer...} & löscht die markierte Ebene. Liegen Objekte darauf, fragt das Programm,
  ob sie auf \gui{Default} wandern (\gui{Move to Default}) oder mitgelöscht werden
  (\gui{Delete objects}); \gui{Cancel} bricht ab. Ein Rückgängig-Schritt; \gui{Default} lässt
  sich nicht löschen. \\
\gui{Make current} & macht die markierte Ebene zur aktuellen: Neue Objekte entstehen dort. \\
Klick auf eine Ebene, dann \gui{Move selection here} & verschiebt die ausgewählten Objekte auf
  die markierte Ebene (ein Rückgängig-Schritt \enquote{Change layer}). \\
\bottomrule
\end{referenz}

\paragraph{Die aktuelle Ebene.} Jedes Dokument hat die Ebene \gui{Default} (hellgrau); ein
neues Dokument beginnt mit ihr als aktueller Ebene. Neu gezeichnete und modellierte, mit
\menu{File,Import STEP...} eingelesene und von Plugins erzeugte Objekte landen auf der
\emph{aktuellen} Ebene (fett in der Liste). Markieren Sie eine Ebene und klicken Sie auf
\gui{Make current}, um sie zur aktuellen zu machen; die Wahl wird mit dem Dokument
gespeichert. Kopien (\cmd{Copy}, Anordnungen) behalten dagegen die Ebene ihres Originals, und
das Ergebnis einer booleschen Operation übernimmt Name, Ebene und Farbe des ersten
ausgewählten Körpers. Wird die aktuelle Ebene gelöscht, ist wieder \gui{Default} aktuell.

\paragraph{Farben.} Ein Objekt wird in der Farbe seiner Ebene gezeichnet, es sei denn, es hat
eine eigene Farbe. Eigene Farben setzen Plugins oder sie stammen aus einer Datei; in der
Oberfläche lassen sie sich nicht ändern. Ausgewählte Objekte erscheinen unabhängig davon gelb.

\paragraph{Reihenfolge.} Die Reihenfolge der Ebenen ist fest; neue Ebenen kommen ans Ende.

\paragraph{Beispiel: Hilfsgeometrie auf eine eigene Ebene.}
\begin{schritte}
  \item Klicken Sie im Dock \gui{Layers} auf \gui{New layer}. Es erscheint \enquote{Layer 1}.
  \item Doppelklicken Sie auf den Namen und tippen Sie \code{Hilfslinien}, dann \keys{Enter}.
  \item Doppelklicken Sie auf das Farbfeld und wählen Sie ein dunkles Grau.
  \item Wählen Sie in den Ansichten die Hilfskurven aus (z.\,B.\ mit einem Kreuzen-Rechteck
    und \keys{Strg}+Klick auf das, was nicht dazugehört).
  \item Klicken Sie im Dock auf die Zeile \enquote{Hilfslinien} und dann auf \gui{Move selection
    here}.
  \item Mit dem Kontrollkästchen vor \enquote{Hilfslinien} blenden Sie die Hilfskurven nun nach
    Bedarf aus und ein.
\end{schritte}

\begin{hinweis}
\gui{Move selection here} verschiebt immer ganze Objekte. Ist nur eine Fläche oder Kante einer
Hülle ausgewählt, wandert die ganze Hülle auf die Ebene.
\end{hinweis}

\begin{hinweis}
Anlegen, Umbenennen, Umfärben und Löschen von Ebenen und das Verschieben von Objekten auf eine
Ebene sind Rückgängig-Schritte. Ein-/Ausblenden, Sperren und die Wahl der aktuellen Ebene sind
es nicht -- sie werden aber mit dem Dokument gespeichert: Das Dokument gilt danach als geändert,
und beim Schließen fragt das Programm nach dem Speichern.
\end{hinweis}

\section{Gruppen}\label{sec:gruppen}
\index{Gruppe}

Gruppen fassen Objekte zu einem Knoten im Modellbaum zusammen. Gruppen können andere Gruppen
enthalten. Sie entstehen auf drei Wegen:

\begin{itemize}
  \item von Hand mit \menu{Edit,Group...} (\keys{Strg+G}) oder \gui{Group selection...} im
    Kontextmenü des Modellbaums,
  \item beim Einlesen einer STEP-Datei: Alle Körper der Datei bilden eine Gruppe mit dem
    Dateinamen ohne Endung (\autoref{sec:step-import}),
  \item durch Plugins: Der Raytracer etwa bildet die Struktur des optischen Systems als
    Gruppen nach (\autoref{kap:raytracer}).
\end{itemize}

\subsection{Gruppe bilden}

\begin{schritte}
  \item Wählen Sie die Objekte aus, die zusammengehören -- in den Ansichten oder im
    Modellbaum.
  \item Drücken Sie \keys{Strg+G} (oder \menu{Edit,Group...}).
  \item Der Dialog \gui{Group} fragt nach dem \gui{Name of the new group:}. Vorgeschlagen ist
    \enquote{Group \textit{n}}. Geben Sie einen Namen ein und bestätigen Sie mit \gui{OK};
    \gui{Cancel} bricht ab. Ein leerer Name ergibt \enquote{Group}.
\end{schritte}

Die Objekte verlassen dabei ihre bisherige Gruppe und kommen in die neue. Lagen \emph{alle}
ausgewählten Objekte vorher in derselben Gruppe, wird die neue Gruppe eine Untergruppe davon;
sonst entsteht sie auf oberster Ebene des Modellbaums. Das Bilden einer Gruppe ist ein
Rückgängig-Schritt (\enquote{Group}).

\begin{tipp}
So bauen Sie eine Hierarchie von außen nach innen auf: erst alle Teile des Objektivs zu
\enquote{Objektiv} gruppieren, dann die beiden Linsen des Kittglieds -- die jetzt beide in
\enquote{Objektiv} liegen -- zu \enquote{Kittglied}. \enquote{Kittglied} wird automatisch
Untergruppe von \enquote{Objektiv}.
\end{tipp}

\subsection{Gruppe auflösen}

\begin{schritte}
  \item Klicken Sie im Modellbaum auf die Zeile der Gruppe. Dabei werden alle ihre Objekte
    ausgewählt.
  \item Klicken Sie mit der rechten Maustaste auf dieselbe Gruppenzeile und wählen Sie
    \gui{Ungroup} -- oder drücken Sie \keys{Strg+Shift+G}.
\end{schritte}

Die Gruppe verschwindet; ihre Objekte und Untergruppen rücken eine Stufe nach oben in die
Gruppe, in der sie lag (oder auf die oberste Ebene). Die Objekte selbst bleiben unverändert.
Auch das ist ein Rückgängig-Schritt (\enquote{Ungroup}). Sind mehrere Gruppenzeilen markiert,
werden alle aufgelöst.

\begin{hinweis}
\gui{Ungroup} wirkt auf die im Modellbaum \emph{markierten Gruppenzeilen}, nicht auf die
Auswahl in den Ansichten. Geschieht nichts, ist keine Gruppenzeile markiert: Klicken Sie die
Gruppenzeile noch einmal an und wiederholen Sie den Befehl.
\end{hinweis}

\subsection{Umbenennen und Sichtbarkeit}

\begin{itemize}
  \item \textbf{Umbenennen:} Doppelklick auf den Gruppennamen im Modellbaum, neuen Namen
    tippen, \keys{Enter}. Das Umbenennen ist ein Rückgängig-Schritt (\enquote{Rename group}).
  \item \textbf{Ein- und ausblenden:} Kontrollkästchen vor der Gruppe. Eine ausgeblendete
    Gruppe blendet alle Objekte und Untergruppen darin aus, ohne deren eigene Kästchen zu
    ändern; eingeschaltet erscheint alles wieder so, wie es vorher war. \gui{Show all} und
    \gui{Isolate} schalten alle Gruppen ein (\autoref{sec:ausblenden}).
\end{itemize}

Objekte selbst lassen sich im Programm nicht umbenennen; ihre Namen stammen aus STEP-Dateien
oder von Plugins.

\section{Der Modellbaum}\label{sec:modellbaum}
\index{Modellbaum}

Das Dock \gui{Model} (\menu{View,Model}) zeigt das Dokument als Baum mit zwei Spalten:

\begin{description}
  \item[Model] -- der Name. Gruppen haben ein Ordnersymbol, Objekte ohne eigenen Namen
    heißen nach Art und Nummer, z.\,B.\ \enquote{Curve 17}. Vor jedem Eintrag steht ein
    Kontrollkästchen für die Sichtbarkeit.
  \item[Type] -- \gui{Group} oder die Art des Objekts: \gui{Point}, \gui{Curve},
    \gui{Surface}, \gui{Trimmed surface}, \gui{Shell}.
\end{description}

Gruppen auf oberster Ebene sind aufgeklappt, tiefere zugeklappt; der Pfeil vor einer Gruppe
klappt sie auf und zu. Objekte mit dem Attribut \code{sim.inactive} erscheinen ausgegraut; ihr
Tooltip lautet \enquote{excluded from computations (sim.inactive)} (\autoref{sec:attribute}).

\paragraph{Auswahl in beide Richtungen.} Ein Klick auf ein Objekt im Baum wählt es in den
Ansichten aus; ein Klick auf eine Gruppe wählt alle Objekte darin aus, auch die in
Untergruppen. Mit \keys{Shift}+Klick markieren Sie einen Bereich von Zeilen, mit
\keys{Strg}+Klick einzelne Zeilen dazu. Umgekehrt markiert der Baum die Zeilen der Objekte,
die Sie in den Ansichten auswählen. Objekte auf einer gesperrten Ebene lassen sich auch hier
nicht auswählen. Das ist der schnellste Weg, ein Objekt in einer großen
Baugruppe zu finden: Anklicken im Baum, dann \keys{Z}, \keys{S} (\gui{Zoom selected}).

\paragraph{Kontextmenü.} Ein Rechtsklick in den Baum öffnet:

\begin{referenz}{>{\raggedright\arraybackslash}p{3.6cm}X}
\toprule
\textbf{Eintrag} & \textbf{Wirkung} \\
\midrule
\endhead
\gui{Group selection...} & wie \menu{Edit,Group...}: die ausgewählten Objekte zu einer neuen
  Gruppe. \\
\gui{Ungroup} & löst die markierten Gruppen auf. \\
\gui{Hide} & blendet die ausgewählten Objekte aus (wie \keys{H}). \\
\gui{Isolate} & zeigt nur die ausgewählten Objekte (wie \menu{Edit,Isolate selected}). \\
\gui{Unhide} & zeigt die einzeln ausgeblendeten Objekte wieder (wie \keys{Shift+H}). \\
\gui{Show all} & schaltet alle Objekte, Gruppen und Ebenen ein. \\
\bottomrule
\end{referenz}

\begin{hinweis}
Das Kontrollkästchen eines \emph{Objekts} im Baum und \gui{Hide}/\gui{Unhide} schalten
dieselbe Eigenschaft. Ein Objekt, dessen Kästchen gesetzt ist, kann trotzdem unsichtbar sein,
wenn seine Gruppe oder seine Ebene ausgeschaltet ist.
\end{hinweis}

\section{Das Dock Properties}\label{sec:eigenschaften}
\index{Eigenschaften}

Das Dock \gui{Properties} (\menu{View,Properties}) zeigt die Eigenschaften der Auswahl. Oben
steht eine Überschrift, darunter eine Tabelle mit den Spalten \gui{Property} und \gui{Value}
und ganz unten der Knopf \gui{Add attribute...}.

\begin{description}
  \item[Nichts ausgewählt:] \enquote{Nothing selected}.
  \item[Ein Objekt oder Teilobjekt:] Die Überschrift nennt den Namen (oder \enquote{Entity
    \textit{Id}}), bei einem Teilobjekt mit Zusatz, z.\,B.\ \enquote{part\_0\_1 $\cdot$ face 3}.
    Die Tabelle beginnt immer mit \gui{Layer} (Name der Ebene) und \gui{Id} (die eindeutige
    Nummer des Objekts), gefolgt von den geometrischen Eigenschaften (siehe unten) und, bei
    ganzen Objekten, den Attributen.
  \item[Mehrere:] Die Überschrift lautet \enquote{\textit{n} items selected}; die Tabelle
    zählt, wie viele Objekte bzw.\ Teilobjekte welcher Art ausgewählt sind, z.\,B.\
    \enquote{Shell 3}, \enquote{face 2}.
\end{description}

Lange Werte werden in der Mitte gekürzt (\dots); der Tooltip zeigt den vollständigen Wert.
Zahlen erscheinen mit bis zu acht gültigen Stellen, Längen in Millimetern, Flächen in
\si{\milli\metre\squared}, Volumina in \si{\milli\metre\cubed}.

\subsection{Geometrische Eigenschaften}

\begin{referenz}{>{\raggedright\arraybackslash}p{3.4cm}X}
\toprule
\textbf{Art (\gui{Type})} & \textbf{Angezeigte Eigenschaften} \\
\midrule
\endhead
\gui{Point} & \gui{Position} ($x, y, z$). \\
\gui{NURBS curve} & \gui{Degree}, \gui{Control points} (Anzahl), \gui{Rational}
  (\code{yes}, wenn nicht alle Gewichte 1 sind), \gui{Knots}, \gui{Domain} (Parameterbereich),
  \gui{Start}, \gui{End}, \gui{Closed}, \gui{Length} (Bogenlänge). \\
\gui{NURBS surface} & \gui{Degree (u, v)}, \gui{Control net} ($n_u \times n_v$),
  \gui{Rational}, \gui{Knots u}, \gui{Knots v}, \gui{Area}. \\
\gui{Trimmed surface} & wie die Fläche, dazu \gui{Trimmed} (\code{yes}, wenn die Fläche
  beschnitten ist), \gui{Holes} (Zahl der Löcher), gegebenenfalls \gui{Label} und
  \gui{Area} (Inhalt der beschnittenen Fläche). \\
\gui{Shell (B-rep)} & \gui{Faces}, \gui{Edges}, \gui{Vertices} (Anzahlen), \gui{Closed}
  (\code{yes} bei einem geschlossenen Körper), \gui{Area} (Oberfläche) und, nur bei
  geschlossenen Hüllen, \gui{Volume}. \\
\gui{Face \textit{n} of a shell} & wie \gui{Trimmed surface}. \\
\gui{Edge \textit{n} of a shell} & \gui{Used by faces} (an wie vielen Flächen die Kante
  anliegt: 2 bei einer inneren Kante, 1 an einem offenen Rand) und die Daten der Kantenkurve
  wie bei \gui{NURBS curve}. \\
\gui{Vertex \textit{n} of a shell} & \gui{Position}. \\
\bottomrule
\end{referenz}

\paragraph{Knotenvektoren} werden kompakt angezeigt: jeder verschiedene Knotenwert einmal,
mehrfache Knoten mit ihrer Vielfachheit, z.\,B.\ \enquote{0 (\texttimes4), 0.5, 1
(\texttimes4)} für einen kubischen Spline mit einem inneren Knoten. Bei mehr als zwölf
verschiedenen Werten stehen nur die ersten und letzten fünf, dazwischen \enquote{\dots}.

\paragraph{Nicht berechenbare Werte} (etwa die Länge einer entarteten Kurve) erscheinen als
\enquote{n/a} mit dem Grund in Klammern.

\begin{tipp}
Die Werte im Dock \gui{Properties} sind exakt aus der NURBS-Geometrie berechnet, nicht aus
der Darstellung. Für Abstände und Winkel zwischen Punkten gibt es die Befehle \cmd{Distance}
und \cmd{Angle}, für Flächeninhalt, Volumen und Masseneigenschaften die Befehle der Gruppe
\gui{Analyze} (\autoref{kap:modellieren}).
\end{tipp}

\section{Freie Attribute}\label{sec:attribute}
\index{Attribut}

Attribute sind Schlüssel-Wert-Paare, die Sie an ein Objekt hängen und die mit dem Dokument
gespeichert werden. Sie haben für das CAD-Programm keine Bedeutung, wohl aber für Plugins: Ein
Plugin kann Objekte an Attributen erkennen, sie überspringen oder seine eigenen Daten daran
speichern. Plugins verwenden dafür Schlüssel mit einem Präfix, etwa \code{rtt.} (Raytracer),
\code{pt.} (Path Tracer), \code{al.} (Automotive Lighting) oder \code{fg.} (FourierGrating).

\paragraph{Anzeigen.} Hat ein ausgewähltes Objekt Attribute, stehen sie im Dock
\gui{Properties} unter der fett gesetzten Zeile \gui{Attributes}, je Zeile Schlüssel und
Wert.

\paragraph{Anlegen.}
\begin{schritte}
  \item Wählen Sie genau \emph{ein} ganzes Objekt aus (kein Teilobjekt). Nur dann ist
    \gui{Add attribute...} anklickbar.
  \item Klicken Sie auf \gui{Add attribute...}. Der Dialog \gui{Add attribute} fragt nach
    \gui{key = value}; vorbelegt ist \code{sim.inactive = true}.
  \item Geben Sie Schlüssel und Wert mit Gleichheitszeichen ein, z.\,B.\
    \code{rtt.role = mirror}, und bestätigen Sie mit \gui{OK}.
\end{schritte}

Der Wert wird nach seiner Form gespeichert: \code{true} oder \code{false} (Groß-/Kleinschreibung
egal) als Wahrheitswert, eine Zahl als Zahl, alles andere als Text. Ohne Gleichheitszeichen
wird der eingegebene Text zum Schlüssel mit dem Wert \code{true}. Ein vorhandener Schlüssel
wird überschrieben.

\paragraph{Ändern und löschen.} Doppelklicken Sie in der Attributzeile auf den Wert, ändern Sie
ihn und bestätigen Sie mit \keys{Enter}. Ein \emph{leerer} Wert löscht das Attribut. Anlegen,
Ändern und Löschen sind Rückgängig-Schritte (\enquote{Change attribute}).

\paragraph{Das Attribut sim.inactive.}\index{sim.inactive} Mit \code{sim.inactive = true}
schließen Sie ein Objekt von Berechnungen aus, ohne es zu löschen oder auszublenden -- etwa
eine Halterung, die im Modell bleiben, aber kein Licht abschatten soll. Das Objekt erscheint im
Modellbaum ausgegraut. Ob und wie ein Plugin das Attribut auswertet, steht in seinem Kapitel:
Der Path Tracer etwa lässt solche Objekte ganz aus der Szene (\autoref{kap:pathtracer}).
Entfernen Sie das Attribut (leerer Wert), nimmt das Objekt wieder teil.

\begin{warnung}
Ändern Sie Attribute, die ein Plugin selbst geschrieben hat (z.\,B.\ \code{rtt.revision}),
nur, wenn das Kapitel des Plugins es beschreibt. Das Plugin verlässt sich auf ihren Inhalt.
\end{warnung}

\section{Rückgängig machen und Wiederherstellen}\label{sec:undo}
\index{Rückgängig}\index{Undo}

\menu{Edit,Undo} (\keys{Strg+Z}) macht den letzten Bearbeitungsschritt rückgängig,
\menu{Edit,Redo} (\keys{Strg+Shift+Z}, unter Windows auch \keys{Strg+Y}) stellt ihn wieder
her. Die Menüeinträge nennen den Schritt, z.\,B.\ \enquote{Undo Move} oder \enquote{Redo
Import STEP}. Die Zahl der Schritte ist nicht begrenzt.

\paragraph{Was ein Schritt ist.}
\begin{itemize}
  \item Jeder Befehl ergibt einen Schritt, auch wenn er mehrere Objekte erzeugt (eine
    Anordnung, eine Polylinie, \dots).
  \item \gui{Delete}, \gui{Group}, \gui{Ungroup}, das Umbenennen einer Gruppe, \gui{Move
    selection here} (\enquote{Change layer}), das Bearbeiten von Attributen (\enquote{Change
    attribute}).
  \item Anlegen, Umbenennen, Umfärben und Löschen von Ebenen (das Löschen mitsamt dem
    Verschieben oder Löschen ihrer Objekte).
  \item Jede Bearbeitung im Node-Editor (\enquote{Edit node workflow}); das Verschieben eines
    Blocks zählt einmal, beim Loslassen.
  \item \menu{File,Import STEP...} -- das Einlesen mitsamt der neuen Gruppe ist ein Schritt.
  \item \gui{Convert to CAD curves...} im Dock \gui{Results} (\enquote{Rays to curves}).
  \item Alles, was \emph{eine} Plugin-Aktion am Modell ändert -- neue, geänderte und gelöschte
    Objekte, Gruppen und Attribute, und auch eine Änderung des Workflows durch das Plugin --,
    ist zusammen ein Schritt. Ein \gui{Undo} nach \enquote{Load system} entfernt also das ganze
    geladene System.
\end{itemize}

\paragraph{Was kein Schritt ist} und daher nicht rückgängig gemacht werden kann:
\begin{itemize}
  \item Ansicht und Zoom, Darstellungsmodus, Auswahl,
  \item Ein- und Ausblenden von Objekten, Gruppen und Ebenen, Sperren von Ebenen, die Wahl der
    aktuellen Ebene (sie werden aber mit dem Dokument gespeichert und markieren es als
    geändert),
  \item Parameter im Dock \gui{Plugin} (sie werden aber mit dem Dokument gespeichert),
  \item Ergebnisse im Dock \gui{Results} (Ein-/Ausblenden, \gui{Delete run}).
\end{itemize}

\paragraph{Grenzen.} \menu{File,New} und \menu{File,Open...} leeren die Liste der Schritte;
nach dem Öffnen kann man also nichts \enquote{vor} das Öffnen zurücknehmen. Speichern leert
sie nicht: Sie können auch nach dem Speichern zurückgehen. Das Dokument gilt dann wieder als
geändert, und der Fenstertitel zeigt das Sternchen.

